% Rapid Communication for Frontiers of Mathematical Finance.
% Constraint (FMF Rapid Communications): main text <= 5 pages; figures,
% tables, and references <= 4 additional pages. Text runs pages 1-5 with no
% floats; bibliography, tables, and the figure follow a \clearpage.
% Build: tectonic paper.tex (from this directory).
\documentclass[11pt]{article}

\input{../shared/preamble.tex}

\usepackage[margin=1in]{geometry}
\ExecuteBibliographyOptions{useprefix=true}
\usepackage{cleveref}

\newtheorem{theorem}{Theorem}
\newtheorem{proposition}{Proposition}
\newtheorem{remark}{Remark}

\newcommand{\E}{\mathbb{E}}
\newcommand{\Prob}{\mathbb{P}}
\newcommand{\R}{\mathbb{R}}
\DeclareMathOperator*{\argminop}{arg\,min}
\DeclareMathOperator*{\argmaxop}{arg\,max}

\title{Thurstone Portfolios:\\ Long-Only Allocation by Inverting Winning Probabilities}
\author{Peter Cotton\thanks{%
  [Affiliation and postal address to be supplied.] Email:
  \texttt{peter.cotton@gsmc.ai}. Code: the \texttt{allocation} and
  \texttt{thurstone} packages, \texttt{github.com/microprediction}.}}
\date{\today}

\begin{document}
\maketitle

\begin{abstract}
We construct long-only portfolios by reading weights as the winning
probabilities of a Thurstonian race among assets: each asset draws a noisy
performance under an arbitrary centered joint law, and its weight is the
probability that it comes out best. Latent abilities are backed out of any
benchmark by a fast exact inverse, and the portfolio is the race re-run under
the estimated dependence. The organizing property is \emph{redundancy
consistency}---a near-duplicate asset splits a position rather than doubling
it---which capitalization and equal weighting violate by axiom (the
red-bus/blue-bus failure of Luce's rule) and which mean--variance attains only
through an unstable inverse of a near-singular covariance. Because the race
samples from a correlation rather than inverting one, the same near-singularity
that breaks $\Sigma^{-1}$ is priced smoothly: the construction is feasible on
the simplex without optimization, reproduces its benchmark exactly, admits
tail-dependent simulations no covariance functional can express, solves an
implied mean--(convex-regularizer) program, and is locally Lipschitz in the
correlation, bounding turnover. We state and prove these properties and
document them empirically: a fourteen-year Dow walk-forward at four to nine
times less turnover than clustering allocators, and a $500$-asset known-truth
market where sample minimum variance degrades from $1.4\times$ to $91\times$
the true minimum variance as $n/T \to 1$ while the race holds flat at
$3.3\times$, long-only, with no inverse taken.
\end{abstract}

\noindent\textbf{Keywords:} portfolio construction; discrete choice; Thurstone
model; independence of irrelevant alternatives; redundancy; tail dependence;
turnover.\\
\textbf{JEL:} G11, C25, C15, C58. \quad \textbf{MSC2020:} 91G10, 62P05, 60E15.

\section{Introduction}
\label{sec:intro}

A classical test of a probabilistic choice model is duplication: split one
alternative into two near-identical copies and ask what happens to the choice
probabilities. Luce's axiom---independence of irrelevant alternatives
(IIA)---forces the copies to draw fresh probability from the whole field, the
red-bus/blue-bus paradox that unseated the logit model in discrete choice
\parencite{mcfadden1974}. Portfolio construction faces the same test daily.
Universes contain near-duplicates---share classes, an ETF beside its dominant
holding, two miners of one commodity---and every weighting rule implicitly
decides how much a near-copy of something already held is worth. Rules that fix
the ratio of two weights independently of the rest of the field
(capitalization weighting on a restricted universe, equal weighting, any
softmax score rule) are \emph{Lucian}: they let duplicated exposure count
twice. Mean--variance de-duplicates correctly in population, but through
$\Sigma^{-1}$, and the direction in which two assets become identical is
exactly the direction in which $\Sigma$ loses rank: the sample answer divides
by a vanishing eigenvalue and swings between the twins. Inversion-free
clustering allocators (hierarchical risk parity, \textcite{lopezdeprado2016};
Schur-complementary allocation, \textcite{cotton2024schur}) de-duplicate by
grouping, but discretely, the split jumping as cluster membership flips.

This paper develops the Thurstonian alternative: treat allocation as a
\emph{choice} among assets, model the choice with Thurstone's law of
comparative judgment \parencite{thurstone1927}, and use the resulting winning
probabilities as weights. Two assets that move together compete with each
other rather than independently, so a co-moving pair collapses smoothly toward
a single position---redundancy consistency without an inverse, with a closed
characterization (\Cref{prop:redundancy,prop:monotone}) and a Lipschitz bound
on turnover (\Cref{thm:smooth}). The Thurstonian (probit) model is older and
strictly more general than the logit; it lost the applied field only on
tractability, the choice probabilities being Gaussian orthant integrals with
an inverse long considered prohibitive. Two developments remove the
obstruction. The \emph{Fast Ability Transform} \parencite{cotton2021}
evaluates the forward and inverse maps of the independent race in linear time
on a lattice, exactly and Monte-Carlo-free; and the correlated forward race is
cheap to simulate, with a common-seed construction that makes the simulated
weights move smoothly with the correlation.

Three properties usually treated separately are then one. Redundancy is
near-singularity, so an allocator that prices redundancy \emph{smoothly} is by
the same token robust to ill-conditioning (it samples from a correlation, never
inverting it) and low in turnover (its weights move only as far as the
correlation moves). A fourth property stands apart: driven by a tail-dependent
simulation, the race responds to lower-tail co-movement that no covariance
functional can encode (\Cref{prop:tail})---the winner of a race is an order
statistic of the joint draw, and an order statistic reads the tail copula
where a second moment cannot. Over 2010--2024, days on which all $28$
continuously listed Dow constituents fell by more than $2\%$ occurred roughly
$80$ times more often than a Gaussian fit to their full mean and covariance
allows; a cluster that crashes together, but only in the crash, is invisible
to every covariance-based allocator and is priced by the race.

We claim no out-of-sample outperformance. The contribution is a construction
with provable structure: feasibility without optimization, exact benchmark
reproduction, redundancy, monotone de-duplication, tail consistency, an
implied convex program, and a turnover bound---together with empirical
evidence that the properties bind in practice (\Cref{sec:evidence}).

\section{The construction}
\label{sec:method}

Each of $n$ assets has a latent ability $a_i \in \R$ and a noisy
\emph{performance}
\begin{equation}
  X_i = a_i + \varepsilon_i, \qquad
  \varepsilon = (\varepsilon_1,\dots,\varepsilon_n) \sim S,
  \label{eq:performance}
\end{equation}
where $S$ is any centered joint law (the \emph{simulation}), and the
\emph{minimum} performance wins, so a smaller ability denotes a stronger
competitor. Writing $\theta = -a$ and
\begin{equation}
  G_S(\theta) = \E_{\eta\sim S}\Big[\max_i\,(\theta_i + \eta_i)\Big], \qquad
  W_S(\theta) = \nabla G_S(\theta) \in \Delta,
  \label{eq:popmap}
\end{equation}
$W_S(\theta)_i = \Prob(X_i = \min_k X_k)$ is the vector of winning
probabilities (the gradient identity is part of \Cref{thm:objective}). The
portfolio uses two laws: a \emph{reference} $S_0$ under which abilities are
\emph{calibrated} to reproduce a target allocation
$w^{\mathrm{target}} \in \operatorname{int}\Delta$ (a capitalization index, an
equal-weight book, any benchmark), and a \emph{target} law $S_1$ carrying the
estimated dependence, under which the race is re-run:
\begin{equation}
  T_{S_0\to S_1}(w^{\mathrm{target}})
  \;=\; W_{S_1}\!\big(W_{S_0}^{-1}(w^{\mathrm{target}})\big).
  \label{eq:tiltmap}
\end{equation}

\begin{proposition}[Calibration]
\label{prop:calib}
For a Gaussian reference $S_0 = \mathcal N(0, C_0)$ with $C_0 \succ 0$, the map
$W_{S_0} : \R^n/\operatorname{span}\{\mathbf 1\} \to \operatorname{int}\Delta$
is a smooth bijection with smooth inverse: every strictly positive benchmark
has a unique ability vector, modulo an additive constant.
\end{proposition}
\begin{proof}
$G_{S_0}$ is convex and, for $C_0 \succ 0$, $C^\infty$, with Hessian
$\nabla^2 G_{S_0} = \nabla_\theta W_{S_0}$ positive definite on
$\{v : \mathbf 1^\top v = 0\}$ and annihilating $\mathbf 1$ (since
$G_{S_0}(\theta + c\mathbf 1) = G_{S_0}(\theta) + c$). Hence $G_{S_0}$ is
strictly convex on the quotient, so its gradient is a diffeomorphism onto the
relative interior of its range, which is $\operatorname{int}\Delta$ since any
coordinate is driven to $0$ or $1$ by sending abilities to $\mp\infty$.
\end{proof}

In the Gaussian workhorse the two laws are two correlation matrices, a
reference $C_{\mathrm{calib}}$ (identity, or a one-factor market model) and
the estimated $C_{\mathrm{tilt}}$, blended by a scalar $\phi \in [0,1]$,
$C_\phi = (1-\phi)\,C_{\mathrm{calib}} + \phi\,\widehat C$, so $\phi = 0$
reproduces the benchmark (\Cref{prop:fixedpoint}) and $\phi = 1$ uses the full
estimate. But nothing restricts $S_1$ to Gaussian: the implementation treats it
as a black-box sampler, so a copula with genuine tail dependence, a
multivariate $t$, or an arbitrary structured generator are equally admissible.

\emph{Computation.} Calibration under the independent reference is the exact
linear-time lattice inverse of \textcite{cotton2021}; under the one-factor
reference the assets are conditionally independent given the factor, so the
same lattice race is integrated over a few Gauss--Hermite nodes with a damped
fixed point. The tilt is Monte Carlo: win frequencies of $M \sim 10^4$--$10^5$
draws, with three provisions that matter in practice. (i) \emph{Common seeds:}
a fixed standard-normal ensemble $Z_1,\dots,Z_M$ is recoloured by
$C_\phi^{1/2}$, so identical correlations give identical weights exactly, and
turnover tracks correlation change rather than sampling noise. (ii) \emph{No
inversion:} only the square root is formed, and only ever multiplied by;
$\operatorname{cond}(C^{1/2}) = \sqrt{\operatorname{cond}(C)}$, and a
rank-deficient correlation merely yields a valid degenerate race. (iii)
\emph{Scale:} a $k$-factor decomposition $C = BB^\top + D$ keeps the transport
$O(Mnk)$, streaming, and well-posed to thousands of names.

\section{Properties}
\label{sec:properties}

Fix the calibrated abilities and view the weights as a function of the target
law.

\begin{proposition}[Feasibility; benchmark reproduction]
\label{prop:fixedpoint}
For every centered $S$ the weights lie in the simplex, $w_i \ge 0$,
$\sum_i w_i = 1$; and if $S_1 = S_0$ then $w = w^{\mathrm{target}}$.
\end{proposition}
\begin{proof}
Each realization awards one unit of weight, split evenly among the assets
attaining the minimum; summing preserves nonnegativity and total mass.
Reproduction is \eqref{eq:tiltmap} at $S_1 = S_0$.
\end{proof}
Long-only, fully invested portfolios therefore come free---no constraints,
penalties, or projection---and the tilt is a controlled perturbation of its
benchmark, departing only as the dependence estimate departs from the
reference.

\begin{proposition}[Redundancy consistency]
\label{prop:redundancy}
Replace asset $k$ by two duplicates $k', k''$ (equal ability, unit variance,
mutual correlation $1$, the same correlation as $k$ with every other asset).
Then, splitting ties evenly,
$w_{k'} = w_{k''} = \tfrac12 w_k^{\setminus}$ and $w_j = w_j^{\setminus}$ for
$j \ne k', k''$, where $w^{\setminus}$ is the allocation of the original field.
\end{proposition}
\begin{proof}
Perfectly correlated unit-variance Gaussians with equal mean are almost surely
equal, so the field minimum, and hence every other asset's winning event, is
unchanged, while the pair ties exactly when $k$ would have won; the even split
halves that probability. (Mutual correlation $1$ is singular, so the statement
is read in the positive-semidefinite limit, by continuity as the correlation
$\uparrow 1$.)
\end{proof}

\begin{proposition}[Diversification monotonicity]
\label{prop:monotone}
Let $T$ be a subgroup of equal-ability assets, independent of the rest of the
field, with common pairwise correlation $\rho$ within $T$. The subgroup's
combined weight is non-increasing in $\rho$, falling from its independent
value at $\rho = 0$ to a single asset's weight at $\rho = 1$.
\end{proposition}
\begin{proof}
Write $X_i = a + \sqrt{\rho}\,Z + \sqrt{1-\rho}\,\varepsilon_i$ for $i \in T$.
By Slepian's inequality \parencite{slepian1962}, raising the off-diagonal
correlations of a Gaussian vector at fixed marginals raises
$\Prob(\min_{i\in T} X_i > t)$ for every $t$, so $\min_{i\in T} X_i$ is
stochastically increasing in $\rho$. The subgroup wins exactly when its
minimum beats the $\rho$-independent minimum of the complement, a probability
that is therefore non-increasing; the endpoints are $|T|$ independent chances
and \Cref{prop:redundancy} respectively.
\end{proof}

At $\rho = 0$ a duplicate is a genuinely fresh bet and the pair's combined
weight properly exceeds $w_k^{\setminus}$; as $\rho \uparrow 1$ it descends
continuously to redundancy coherence (\Cref{tab:redundancy}). A Lucian rule is
pinned at the $\rho$-blind headcount value for every $\rho$.

\begin{proposition}[Tail consistency]
\label{prop:tail}
Take $X_i$ centred loss-side performances (larger is worse, minimum wins). Let
$G$ be a subgroup of equal-ability assets with $M_G = \min_{i\in G} X_i$ and
$Y = \min_{j\notin G} X_j$ independent of $G$, so $w_G = \Prob(M_G < Y)$.
(i) If a change of the subgroup's joint law raises $M_G$ in first-order
stochastic order while fixing $Y$ and the marginals, then $w_G$ does not
increase. (ii) In the comonotone limit $X_i = Z_G$ for all $i \in G$, $w_G$
equals a single competitor's winning probability, independent of $|G|$, split
equally within $G$.
\end{proposition}
\begin{proof}
(i) $w_G = \E[\overline F_Y(M_G)]$ with $\overline F_Y$ non-increasing. (ii)
Comonotonicity gives $M_G = Z_G$, the winning event of one competitor;
exchangeability splits it.
\end{proof}
Deepening joint lower-tail co-movement at fixed marginals---and even at fixed
correlation---raises $M_G$ stochastically, so a cluster that crashes together
is de-weighted toward a single bet even when its covariance is unchanged: the
effective number of bets is state-dependent, $|G|$ in calm states and $1$ in
the crash, a collapse a covariance averages away. Driving the race with a
downside (lower-partial-moment) covariance recovers most of the effect; a
tail-dependent simulation is exact.

\begin{theorem}[Implied objective]
\label{thm:objective}
For any centered $S$, $G_S$ in \eqref{eq:popmap} is convex; the tie-split
winning vector is a selected subgradient
$w \in \partial G_S(\theta) = \argmaxop_{p\in\Delta}\{\langle\theta,p\rangle -
\Omega_S(p)\}$ with $\Omega_S = G_S^{*}$ convex on the simplex, and when $S$ is
full-rank continuous, $w = \nabla G_S(\theta)$ uniquely.
\end{theorem}
\begin{proof}
$G_S$ is an expectation of affine functions of $\theta$, hence convex; with an
almost-surely unique maximizer the envelope theorem gives
$\partial G_S/\partial\theta_i = \Prob(i = \argmaxop_k(\theta_k+\eta_k)) = w_i$
(the Williams--Daly--Zachary identity;
\cite{williams1977,mcfadden1981}). The variational form is Fenchel--Young
duality.
\end{proof}
The portfolio thus maximizes ability-weighted exposure minus a convex
\emph{diversification penalty} determined by the dependence law---the
perturbed-optimizer construction \parencite{berthet2020}. The perturbation is
decisive: Gumbel noise yields Shannon entropy and softmax, a
correlation-\emph{blind} Lucian allocator; Gaussian noise with correlation $C$
makes $\Omega_C$ correlation-aware, which is exactly
\Cref{prop:redundancy,prop:monotone}; a tail-dependent law makes it tail-aware
(\Cref{prop:tail}). Calibration is then reverse optimization in the sense of
\textcite{black1992}: the benchmark implies the abilities under the reference
geometry, and the tilt expresses a view on \emph{dependence}---including a
tail view Black--Litterman cannot state---while the abilities, hence the
return-bearing content of the benchmark, are never touched.

\begin{theorem}[Smoothness and turnover]
\label{thm:smooth}
For fixed abilities, on any compact set of correlation matrices with smallest
eigenvalue bounded below by $\lambda > 0$, the Gaussian-race weight map is
$C^\infty$ with
$\|w(C_1) - w(C_0)\|_1 \le L_\lambda\,\|C_1 - C_0\|_F$, where
$L_\lambda < \infty$ and $L_\lambda \to \infty$ as $\lambda \downarrow 0$.
\end{theorem}
\begin{proof}[Proof sketch]
Asset $i$ wins iff $Y \sim \mathcal N(0,C)$ lands in the fixed polyhedral cone
$A_i = \{y : y_i - y_j \le a_j - a_i\ \forall j\}$, so
$w_i(C) = \int_{A_i} \varphi_C$. Smoothness is differentiation under the
integral on the positive-definite cone. For the bound, Price's theorem
\parencite{price1958}, $\partial\varphi_C/\partial C_{jk} =
\partial^2\varphi_C/\partial y_j\partial y_k$ for $j \ne k$, followed by the
divergence theorem over $A_i$, expresses $\partial w_i/\partial C_{jk}$ as an
integral of $\varphi_C$ over the pairwise-tie faces of $A_i$, which is finite
and controlled by the conditioning of $C$. The common-seed realization
inherits the population motion with difference noise
$O(\sqrt{\|\Delta C\|_F/M})$ and is exactly zero at $\Delta C = 0$.
\end{proof}

\section{Empirical evidence}
\label{sec:evidence}

Three exhibits document that the properties bind; none claims risk-adjusted
outperformance.

\emph{Redundancy} (\Cref{tab:redundancy}). Three equal-ability assets, two of
them duplicates at mutual correlation $\rho_{12}$: the pair's combined weight
descends monotonically from the Lucian headcount $2/3$ at $\rho_{12} = 0$ to
the redundancy-coherent $1/2$ at $\rho_{12} = 1$, as
\Cref{prop:redundancy,prop:monotone} require.

\emph{Turnover} (\Cref{tab:dow}). A daily walk-forward over the $28$ Dow
constituents with full 2010--2024 history, all estimators updated online from
a $252$-day warm-up. The tilt trades four to nine times less than minimum
variance and the clustering allocators at comparable gross Sharpe: its weights
move with the correlation (\Cref{thm:smooth}), not with the re-clustering or
re-seriation those methods incur each step. Charging $10$ basis points on
turnover, only the near-static allocators (equal weight, inverse variance,
risk parity) net more.

\emph{High dimension} (\Cref{tab:highdim}). A known-truth market---$500$ names
in $10$ sectors, one market and one sector factor each, so the true
minimum-variance portfolio and its variance are exact---with the estimation
window $T$ shrunk past $n$. Sample minimum variance passes from best-on-table
($1.4\times$ the true minimum at $n/T = 0.25$: when the data supports the
inversion, invert) through $16\times$ at $n/T = 0.95$ and $91\times$ at
$0.99$, to pure rounding noise beyond $n/T = 1$, where the sample covariance
is singular and $\Sigma^{-1}\mathbf 1$ is undefined. The low-rank race
($k = 12$ factors) holds flat at $3.3\times$ throughout---marginally below
equal weight at $3.5\times$---long-only, leverage $1$, with no qualitative
change at the singularity. Its high-dimensional claim is precisely this:
defined, stable, and correlation-aware in the regime where inversion-based
weights are noise, at $O(Mnk)$ per update.

\emph{Tail consistency} (\Cref{fig:tail}) is illustrated with one independent
asset racing a $99$-name cluster whose lower-tail dependence is swept at fixed
margins: the cluster's combined weight falls toward the two-body limit
$\tfrac12$, the full covariance stays tail-blind, and a downside-covariance
proxy recovers most of the gap.

\section{Discussion}
\label{sec:discussion}

Reading weights as winning probabilities yields a long-only allocator that is
feasible without optimization, anchored to any benchmark, inversion-free and
therefore well-posed at any $n/T$, smooth enough for online use, and---alone
among the allocators discussed---responsive to tail dependence. The same
machinery is a general simplex-regression device (winning probabilities as
redundancy-aware credit among collinear regressors), with forecast combination
the natural application. A companion manuscript in preparation develops the
reverse-optimization geometry (the Bregman form of the tilt and its displaced
Markowitz inverse), concentration control, the CAPM restriction argument, and
the forecast-combination evidence at full length.

\clearpage
\printbibliography

\clearpage

\begin{table}[p]
  \centering
  \caption{Redundancy consistency and diversification monotonicity. Three
  equal-ability assets; assets $1,2$ are duplicates with mutual correlation
  $\rho_{12}$, asset $0$ independent of them ($M = 4\times10^5$ paths). The
  duplicated pair's combined weight falls monotonically from the Lucian
  headcount $2/3$ to the redundancy-coherent $1/2$
  (\Cref{prop:redundancy,prop:monotone}).}
  \label{tab:redundancy}
  \begin{tabular}{ccc}
    \toprule
    $\rho_{12}$ & $w_0$ & $w_1 + w_2$ \\
    \midrule
    $0.00$ & $0.333$ & $0.667$ \\
    $0.50$ & $0.384$ & $0.616$ \\
    $0.90$ & $0.449$ & $0.551$ \\
    $0.99$ & $0.484$ & $0.516$ \\
    $1.00$ & $0.500$ & $0.500$ \\
    \bottomrule
  \end{tabular}
\end{table}

\begin{table}[p]
  \centering
  \caption{Walk-forward over $28$ Dow constituents, 2010--2024, daily online
  updates. Turnover is the mean one-way fraction of the book traded per day;
  net Sharpe charges $10$ basis points on it. The tilt turns over four to nine
  times less than minimum variance and the clustering allocators at comparable
  gross Sharpe (\Cref{thm:smooth}); no risk-adjusted-return claim is made.}
  \label{tab:dow}
  \begin{tabular}{lccc}
    \toprule
    Method & Turnover & Sharpe & Net Sharpe \\
    \midrule
    Equal weight              & $0.000$ & $0.74$ & $0.74$ \\
    Inverse variance          & $0.005$ & $0.78$ & $0.76$ \\
    Risk parity               & $0.004$ & $0.78$ & $0.76$ \\
    Minimum variance          & $0.045$ & $0.77$ & $0.60$ \\
    Hierarchical risk parity  & $0.034$ & $0.79$ & $0.67$ \\
    Schur ($\gamma=0.5$)      & $0.071$ & $0.78$ & $0.53$ \\
    Thurstone tilt            & $0.008$ & $0.71$ & $0.69$ \\
    \bottomrule
  \end{tabular}
\end{table}

\begin{table}[p]
  \centering
  \caption{High dimension: known-truth market with $n = 500$ names in $10$
  sectors (one market factor, one sector factor each), one simulated history,
  estimation window $T$ its first $T$ days. Out-of-sample variance is the
  weight vector's variance under the \emph{true} covariance, as a multiple of
  the true global-minimum variance; the true GMV has leverage $3.4$ and equal
  weight sits at $3.5\times$. Sample minimum variance is
  $w \propto S^{-1}\mathbf 1$; beyond $n/T = 1$ the sample covariance is
  singular and the solver output is rounding noise, reported for illustration.
  The race uses a rank-$12$ factor decomposition of the sample correlation and
  $M = 8192$ common-seed paths.}
  \label{tab:highdim}
  \begin{tabular}{cc cc c}
    \toprule
    & & \multicolumn{2}{c}{sample minimum variance} & race ($k=12$) \\
    \cmidrule(lr){3-4}\cmidrule(lr){5-5}
    $T$ & $n/T$ & OOS var & leverage & OOS var \\
    \midrule
    $2000$ & $0.25$ & $1.4\times$ & $4.2$ & $3.3\times$ \\
    $1000$ & $0.50$ & $2.0\times$ & $5.5$ & $3.3\times$ \\
    $750$  & $0.67$ & $2.7\times$ & $6.4$ & $3.3\times$ \\
    $600$  & $0.83$ & $5.8\times$ & $9.7$ & $3.3\times$ \\
    $525$  & $0.95$ & $16\times$  & $16$  & $3.3\times$ \\
    $505$  & $0.99$ & $91\times$  & $40$  & $3.3\times$ \\
    $400$  & $1.25$ & $1.7\times10^4$ (noise) & $523$ & $3.3\times$ \\
    \bottomrule
  \end{tabular}
\end{table}

\begin{figure}[p]
  \centering
  \includegraphics[width=0.80\textwidth]{figures/tail_consistency.png}
  \caption{Tail consistency (\Cref{prop:tail}). One independent asset versus a
  $99$-name cluster with Clayton lower-tail dependence $\lambda_L$ (assets
  crash together, rally independently; standard-normal margins;
  $M = 1.5\times10^5$; loss-side race on centred performances). As
  $\lambda_L \to 1$ within this family the cluster's combined weight falls
  toward the two-body limit $\tfrac12$: the $99$ names de-duplicate to a
  single horse. A downside-semicovariance drive recovers almost all of the
  effect; the full covariance is tail-blind and stays far above.}
  \label{fig:tail}
\end{figure}

\end{document}
